\frfilename{mt221.tex}
\versiondate{2.6.03}
\copyrightdate{1995}

\def\chaptername{Teorema Dasar Kalkulus}
\def\sectionname{Teorema Vitali di $\Bbb R$}

\newsection{221}
\def\headlinesectionname{Teorema Vitali di $\eightBbb R$}

Saya membahas teorema pertama bab ini dalam satu bagian tersendiri.
Teorema tersebut menempati posisi di antara teori ukuran dan geometri
(bahkan, teorema ini merupakan salah satu hasil dasar dalam `teori ukuran
geometrik'), dan pembuktiannya melibatkan ukuran maupun
geometri garis bilangan real.

\leader{221A}{Teorema Vitali} Misalkan $A$ suatu subhimpunan berbatas dari
$\Bbb R$ dan $\Cal I$ suatu keluarga interval tertutup dalam
$\Bbb R$ yang masing-masing terdiri atas lebih dari satu titik, sedemikian
sehingga setiap titik dari $A$ termasuk dalam
anggota-anggota $\Cal I$ dengan panjang sekecil apa pun. Maka terdapat
suatu himpunan terhitung $\Cal I_0\subseteq\Cal I$ sedemikian sehingga (i)
$\Cal I_0$ saling lepas, yakni $I\cap I'=\emptyset$ untuk setiap pasangan
berbeda $I$, $I'\in\Cal I_0$; (ii)
$\mu(A\setminus\bigcup\Cal I_0)=0$, dengan $\mu$ ukuran Lebesgue pada
$\Bbb R$.

\proof{{\bf (a)} Jika terdapat suatu himpunan berhingga yang saling lepas
$\Cal I_0\subseteq\Cal I$ sedemikian sehingga
$A\subseteq\bigcup\Cal I_0$ (termasuk kemungkinan bahwa
$A=\Cal I_0=\emptyset$), kita selesai. Jadi, mulai sekarang andaikan
bahwa tidak ada $\Cal I_0$ semacam itu.

Misalkan $\mu\sp*$ adalah ukuran luar Lebesgue pada $\Bbb R$. Andaikan
bahwa $|x|<M$ untuk setiap $x\in A$, dan tetapkan

\Centerline{$\Cal I'=\{I:I\in\Cal I,\,I\subseteq[-M,M]\}$.}

\medskip

{\bf (b)} Dalam hal ini, jika $\Cal I_0$ adalah sembarang subhimpunan
berhingga dari $\Cal I'$ yang saling lepas, terdapat suatu $J\in\Cal I'$
yang tidak beririsan dengan anggota mana pun dari $\Cal I_0$. \Prf\ Ambil
$x\in A\setminus\bigcup\Cal I_0$. Sekarang terdapat $\delta>0$ sedemikian
sehingga $[x-\delta,x+\delta]$ tidak bertemu dengan anggota mana pun dari
$\Cal I_0$, dan karena $|x|<M$ kita dapat mengandaikan bahwa
$[x-\delta,x+\delta]\subseteq[-M,M]$. Misalkan $J$ suatu anggota
$\Cal I$ yang memuat $x$ dan panjangnya paling besar $\delta$; maka
$J\in\Cal I'$ dan $J\cap\bigcup\Cal I_0=\emptyset$.\
\Qed

\medskip


{\bf (c)} Sekarang kita dapat memilih suatu barisan
$\sequencen{\gamma_n}$ dari bilangan real dan suatu barisan saling lepas
$\sequencen{I_n}$ di dalam $\Cal I'$ secara induktif sebagai berikut.
Diberikan $\langle I_j\rangle_{j<n}$ (jika $n=0$, ini adalah barisan
kosong, tanpa anggota), dengan $I_j\in\Cal I'$ untuk setiap $j<n$, dan
$I_j\cap I_k=\emptyset$ untuk $j<k<n$, tetapkan

\Centerline{$\Cal J_n=\{I:I\in\Cal I',\,I\cap I_j=\emptyset$ untuk setiap
$j<n\}$.}

\noindent Kita tahu dari (b) bahwa $\Cal J_n\ne\emptyset$. Tetapkan

\Centerline{$\gamma_n=\sup\{\mu I:I\in\Cal J_n\}$;}

\noindent maka $0<\gamma_n\le 2M$. Karena itu, kita dapat memilih suatu
$I_n\in\Cal J_n$ sedemikian sehingga $\mu I_n\ge{1\over 2}\gamma_n$, dan
ini melanjutkan induksi.

\medskip

{\bf (e)} Karena himpunan-himpunan $I_n$ saling lepas, terukur Lebesgue,
dan termuat dalam $[-M,M]$, kita mempunyai

\Centerline{$\sum_{n=0}^{\infty}\gamma_n
\le 2\sum_{n=0}^{\infty}\mu I_n\le 4M<\infty$,}

\noindent dan $\lim_{n\to\infty}\gamma_n=0$. Sekarang definisikan
$I'_n$ sebagai interval tertutup yang berbagi titik tengah dengan
$I_n$ tetapi panjangnya lima kali panjang interval semula, sehingga
interval tersebut menjulur melampaui setiap ujung $I_n$ setidak-tidaknya sejauh
$\gamma_n$. Saya menyatakan bahwa, untuk setiap $n$,

\Centerline{$A\subseteq\bigcup_{j<n}{I}_j\cup\bigcup_{j\ge n}I_j'$.}

\noindent\Prf\Quer\ Andaikan, jika mungkin, bahwa pernyataan ini salah.
Ambil sembarang $x$ yang termasuk dalam
$A\setminus(\bigcup_{j<n}{I}_j\cup\bigcup_{j\ge n}I'_j)$. Misalkan
$\delta>0$ sedemikian sehingga


\Centerline{$[x-\delta,x+\delta]
\subseteq[-M,M]\setminus\bigcup_{j<n}I_j$,}

\noindent dan misalkan $J\in\Cal I$ sedemikian sehingga

\Centerline{$x\in J\subseteq[x-\delta,x+\delta]$.}

\noindent Maka

\Centerline{$\mu J>0=\lim_{m\to\infty}\gamma_m$;}

\noindent misalkan $m$ bilangan bulat terkecil yang lebih besar daripada
atau sama dengan $n$ sedemikian sehingga $\gamma_m<\mu J$. Dalam hal ini,
$J$ tidak mungkin termasuk dalam $\Cal J_m$, sehingga pasti terdapat
suatu $k<m$ sedemikian sehingga $J\cap I_k\ne\emptyset$, sebab tentu saja
$J\in\Cal I'$. Berdasarkan pemilihan $\delta$, $k$ tidak mungkin lebih
kecil daripada $n$, sehingga $n\le k<m$, dan $\gamma_k\ge\mu J$. Dalam
hal ini, jarak dari $x$ ke titik ujung $I_k$ yang terdekat paling besar
$\mu J\le \gamma_k$. Namun, ujung-ujung $I'_k$ menjulur melampaui
ujung-ujung $I_k$ setidak-tidaknya sejauh $\gamma_k$, sehingga
$x\in I'_k$; ini bertentangan dengan pemilihan $x$. \Bang \Qed

\medskip

{\bf (f)} Akibatnya,

\Centerline{$\mu\sp*(A\setminus\bigcup_{j<n}I_j)
\le\mu(\bigcup_{j\ge n}I'_j)\le\sum_{j=n}^{\infty}\mu I'_j
\le 5\sum_{j=n}^{\infty}\mu I_j$.}

\noindent Karena

\Centerline{$\sum_{j=0}^{\infty}\mu I_j\le 2M<\infty$,}

\noindent kita harus mempunyai

\Centerline{$\lim_{n\to\infty}\mu\sp*(A\setminus\bigcup_{j<n}I_j)=0,$}

\noindent dan

\Centerline{$\mu(A\setminus\bigcup_{j\in\Bbb N}I_j)
=\mu\sp*(A\setminus\bigcup_{j\in\Bbb N}I_j)
\le\inf_{n\in\Bbb N}\mu^*(A\setminus\bigcup_{j<n}I_j)
=0$.}

Jadi, dalam hal ini kita dapat menetapkan
$\Cal I_0=\{I_n:n\in\Bbb N\}$ untuk memperoleh suatu keluarga terhitung
yang saling lepas di dalam $\Cal I$ dengan
$\mu(A\setminus\bigcup\Cal I_0)=0$.
}%end of proof of 221A

\cmmnt{
\leader{221B}{Catatan (a)} Saya telah menyatakan teorema ini dalam bentuk
`terdapat suatu himpunan terhitung $\Cal I_0\subseteq\Cal I$ sedemikian
sehingga $\ldots$' dalam upaya menemukan cara ringkas untuk menyatakan
tiga kemungkinan

\inset{(i) $A=\Cal I=\emptyset$, sehingga kita harus mengambil
$\Cal I_0=\emptyset$;}

\inset{(ii) terdapat $I_0,\ldots,I_n\in\Cal I$ yang saling lepas
sedemikian sehingga $A\subseteq I_0\cup\ldots\cup I_n$, sehingga kita dapat
mengambil $\Cal I_0=\{I_0,\ldots,I_n\}$;}

\inset{(iii) terdapat suatu barisan saling lepas $\sequencen{I_n}$ di
dalam $\Cal I$ sedemikian sehingga
$\mu(A\setminus\bigcup_{n\in\Bbb N}I_n)=0$, sehingga kita dapat mengambil
$\Cal I_0=\{I_n:n\in\Bbb N\}$.}

\noindent Tentu saja banyak penerapan, seperti pembuktian 221A sendiri,
akan menggunakan bentuk-bentuk dari ketiga alternatif ini.

\header{221Bb}{\bf (b)} Teorema sebagaimana dinyatakan di sini akan
digunakan dalam bagian berikutnya. Namun, prinsip pembuktiannya sama
pentingnya dengan pernyataan teorema tersebut. Interval
$I_n$ dipilih `secara serakah', yakni ketika tiba saatnya memilih $I_n$,
kita melihat keluarga $\Cal J_n$ dari interval-interval yang mungkin,
dengan memperhitungkan pilihan $I_0,\ldots,I_{n-1}$ yang telah dibuat,
dan memilih suatu $I_n\in\Cal J_n$ yang `kira-kira' sebesar mungkin.
Supremum dari kemungkinan-kemungkinan untuk $\mu I_n$ adalah $\gamma_n$;
tetapi karena kita tidak tahu apakah terdapat suatu $I\in\Cal J_n$
sedemikian sehingga $\mu I=\gamma_n$, kita harus puas dengan sesuatu yang
sedikit lebih kecil. Saya mengikuti rumus baku dengan mengambil
$\mu I_n\ge\bover12\gamma_n$, tetapi tentu saja saya dapat mengambil
$\mu I_n\ge\bover{99}{100}\gamma_n$, atau
$\mu I_n\ge(1-2^{-n})\gamma_n$, jika hal itu membantu nantinya.
Hal yang luar biasa ialah bahwa cara ini berhasil; kita dapat memilih
$I_n$ tanpa melihat ke depan dan tanpa meninjau hubungan timbal balik di
antara interval-interval tersebut (bahkan, tanpa memeriksa himpunan $A$)
selain persyaratan minimum bahwa $I_n\cap I_j=\emptyset$ untuk $j<n$, dan
bahkan prosedur yang sembarang dan sambil lalu ini menghasilkan barisan
yang sesuai.

\header{221Bc}{\bf (c)} Saya telah menyatakan teorema ini untuk himpunan
berbatas $A$ dan interval tertutup, yang memadai bagi keperluan kita,
tetapi perubahan sangat kecil pada pembuktiannya cukup untuk menangani
interval sembarang (yang terdiri atas lebih dari satu titik), dan
penyempurnaan lain menangani himpunan tak-berbatas $A$. (Lihat 221Ya.)
}%end of comment

\exercises{
\leader{221X}{Latihan dasar (a)} Misalkan $\alpha\in\ooint{0,1}$.
Andaikan, dalam bagian (c) dari pembuktian 221A, kita mengambil
$\mu I_n\ge\alpha\gamma_n$ untuk setiap $n\in\Bbb N$, alih-alih
$\mu I_n\ge\bover12\gamma_n$. Konstanta berapakah yang tepat untuk
menggantikan $5$ ketika mendefinisikan himpunan-himpunan $I_j'$ dalam
bagian (e)?

\leader{221Y}{Latihan lanjutan (a)}
Misalkan $A$ suatu subhimpunan dari $\Bbb R$ dan $\Cal I$ suatu keluarga
interval dalam $\Bbb R$ yang masing-masing terdiri atas lebih dari satu
titik, sedemikian sehingga setiap titik dari $A$ termasuk dalam
anggota-anggota $\Cal I$ dengan panjang sekecil apa
pun. Tunjukkan bahwa terdapat suatu himpunan terhitung yang saling lepas
$\Cal I_0\subseteq\Cal I$ sedemikian sehingga
$A\setminus\bigcup\Cal I_0$ terabaikan Lebesgue. \Hint{terapkan 221A pada
himpunan $A\cap\ooint{n,n+1}$ dan keluarga
$\{\overline{I}:I\in\Cal I,\,\overline{I}\subseteq\ooint{n,n+1}\}$,
dengan menuliskan $\overline{I}$ untuk interval tertutup yang memiliki
titik-titik ujung yang sama dengan $I$.}

\header{221Yb}{\bf (b)} Misalkan $\Cal J$ sembarang keluarga interval
dalam $\Bbb R$ yang masing-masing terdiri atas lebih dari satu titik.
Tunjukkan bahwa $\bigcup\Cal J$ terukur
Lebesgue.
\Hint{terapkan (a) pada $A=\bigcup\Cal J$ dan keluarga $\Cal I$ yang
terdiri atas subinterval-subinterval dari anggota-anggota $\Cal J$,
dengan setiap subinterval tersebut memuat lebih dari satu titik.}

\header{221Yc}{\bf (c)} Misalkan $(X,\rho)$ suatu ruang metrik, $A$ suatu
subhimpunan dari $X$, dan $\Cal I$ suatu keluarga bola tertutup berjari-jari
positif dalam $X$ sedemikian sehingga setiap titik dari $A$ termasuk dalam
anggota-anggota $\Cal I$ yang sekecil apa pun.
(Di sini saya mengatakan bahwa suatu himpunan adalah `bola tertutup
berjari-jari positif' jika himpunan tersebut dapat dinyatakan dalam bentuk
$B(x,\delta)=\{y:\rho(y,x)\le\delta\}$ dengan $x\in X$ dan $\delta>0$.
Tentu saja, bola semacam itu mungkin saja merupakan himpunan satu titik $\{x\}$.)
Tunjukkan bahwa entah $A$ dapat ditutupi oleh suatu keluarga berhingga
yang saling lepas di dalam $\Cal I$, atau terdapat suatu barisan saling
lepas $\sequencen{B(x_n,\delta_n)}$ di dalam $\Cal I$ sedemikian sehingga

\Centerline{$A\subseteq\bigcup_{m\le n}B(x_m,\delta_m)
  \cup\bigcup_{m>n}B(x_m,5\delta_m)$ untuk setiap $n\in\Bbb N$}

\noindent atau terdapat suatu barisan saling lepas
$\sequencen{B(x_n,\delta_n)}$ di dalam $\Cal I$ sedemikian sehingga
$\inf_{n\in\Bbb N}\delta_n>0$.

\spheader 221Yd Berikan contoh suatu keluarga $\Cal I$ dari interval
terbuka sedemikian sehingga setiap titik dari $\Bbb R$ termasuk dalam
anggota-anggota $\Cal I$ yang sekecil apa pun, tetapi jika
$\sequencen{I_n}$ adalah sembarang barisan saling lepas di dalam $\Cal I$,
dan untuk setiap $n\in\Bbb N$ kita menuliskan $I'_n$ untuk interval
tertutup yang berbagi titik tengah dengan $I_n$ dan panjangnya sepuluh kali
panjang interval semula, maka terdapat suatu $n$ sedemikian sehingga
$\ooint{0,1}\not\subseteq\bigcup_{m<n}I_m\cup\bigcup_{m\ge n}I'_m$.
%mt22bits

\spheader 221Ye(i) Tunjukkan bahwa jika $\Cal I$ adalah suatu keluarga
{\it berhingga} dari interval-interval dalam $\Bbb R$, terdapat
$\Cal I_0$, $\Cal I_1\subseteq\Cal I$ sedemikian sehingga
$\bigcup(\Cal I_0\cup\Cal I_1)=\bigcup\Cal I$ dan $\Cal I_0$ maupun
$\Cal I_1$ merupakan keluarga yang saling lepas. \Hint{lakukan induksi
pada $\#(\Cal I)$.} (ii) Andaikan $\Cal I$ adalah suatu keluarga interval
yang masing-masing terdiri atas lebih dari satu titik dan memiliki panjang
paling besar $1$, yang menutupi suatu himpunan berbatas
$A\subseteq\Bbb R$, dan bahwa $\epsilon>0$. Tunjukkan bahwa
terdapat suatu subkeluarga saling lepas $\Cal I_0$ dari $\Cal I$ sedemikian
sehingga
$\mu^*(A\setminus\bigcup\Cal I_0)\le\bover12\mu^*A+\epsilon$.
\Hint{dengan mengganti setiap anggota $\Cal I$ dengan interval yang sedikit
lebih panjang dan bertitik ujung rasional, reduksikan ke kasus ketika
$\Cal I$ terhitung dan kemudian ke kasus ketika $\Cal I$ berhingga;
sekarang gunakan (i).}
(iii) Gunakan (ii) untuk membuktikan Teorema Vitali. (Saya mempelajari
argumen ini dari J.Aldaz.)
%221A
}%end of exercises

\endnotes{
\Notesheader{221} Saya memberi bagian ini judul `Teorema Vitali di
$\Bbb R$' karena terdapat versi berdimensi $r$, yang akan muncul dalam
Bab 26 di bawah.
Terdapat suatu anomali pada posisi teorema ini. Teorema ini merupakan
unsur yang sangat diperlukan dalam pembuktian beberapa teorema terpenting
dalam teori ukuran; di sisi lain, gagasan-gagasan yang terlibat dalam
pembuktiannya sendiri tidak digunakan di tempat lain dalam teori
elementer. Karena itu, ketika mengajarkan materi ini saya sendiri kadang
melewatkan pembuktiannya, dan saya tidak akan menyalahkan mahasiswa mana
pun yang untuk sementara mengesampingkannya. Tentu saja, pada suatu tahap
setiap ahli teori ukuran harus menguasai metodenya, bukan sekadar demi
kelengkapan, melainkan juga untuk memperoleh intuisi mengenai
variasi-variasi yang mungkin. Saya harus menekankan bahwa {\it prinsip}
pembuktiannya, bukan perinciannya, yang penting, sebab terdapat tak
terbilang banyak bentuk `Teorema Vitali'. (Saya menawarkan beberapa
variasi dalam latihan-latihan di sini dan dalam \S261 di bawah, dan
terdapat banyak variasi lain yang penting dalam kajian lebih lanjut;
salah satunya akan muncul dalam \S472 di Jilid 4.) Saya kira prinsip ini
adalah bahwa

\inset{(i) kita memilih $I_n$ secara serakah, menurut suatu kriteria yang
kurang lebih alami dan berlaku untuk setiap $I_n$ ketika tiba saatnya
memilih interval tersebut, tanpa berusaha melihat ke depan;}

\inset{(ii) kita membuktikan bahwa ukuran interval-interval tersebut menuju nol, meskipun
kita tampaknya tidak melakukan apa pun untuk memastikan bahwa hal itu
akan terjadi (tetapi perhatikan peralihan dari $\Cal I$ ke $\Cal I'$ dalam
bagian (a) dari pembuktian 221A, yang persis diperlukan agar langkah ini
berhasil);}

\inset{(iii) kita memeriksa bahwa untuk definisi $I_n'$ yang sesuai,
dengan memperbesar $I_n$, kita akan mempunyai
$A\subseteq\bigcup_{m<n}I_m\cup\bigcup_{m\ge n}I'_m$
untuk setiap $n$, sedangkan $\sum_{n=0}^{\infty}\mu I'_n<\infty$.}

\noindent Dalam arti tertentu, kita harus menganggap diri kita beruntung
setiap kali cara ini berhasil. Alasan mempelajari sebanyak mungkin variasi
dari teknik semacam ini adalah agar kita belajar menebak kapan kita
mungkin beruntung.
}%end of comment

\discrpage
